\documentclass[11pt]{article}

\usepackage{fullpage}
\usepackage{amsmath}
\usepackage{amssymb}
\usepackage{amsfonts}
\usepackage{mathtools}
\usepackage{enumitem}
\usepackage[colorlinks,linkcolor=magenta,citecolor=blue]{hyperref}

\newcommand{\Rbar}{\bar R}
\newcommand{\Om}{\Omega}
\newcommand{\E}{\mathbb E}
\newcommand{\Var}{\operatorname{Var}}
\newcommand{\tr}{\operatorname{tr}}
\newcommand{\diag}{\operatorname{diag}}
\newcommand{\Unif}{\operatorname{Unif}}
\newcommand{\IG}{\operatorname{IG}}
\newcommand{\Gam}{\operatorname{Gamma}}
\newcommand{\Normal}{\mathcal N}

\title{SuSiE-relax: Prototype and Lite Variational Inference}
\author{}
\date{}

\begin{document}
\maketitle

Computation note of the susie-relax model.

Throughout, \(x\in\mathbb R^J\) is the vector of marginal summary
statistics, \(N\) is the GWAS sample size, and \(\Rbar\) is a reference
LD correlation matrix. We write
\[
\Om=\Rbar^{-1},\qquad
\mu_j=\Rbar_{-j,j},\qquad
S_j=\Rbar_{-j,-j}-\Rbar_{-j,j}\Rbar_{j,-j}.
\]
For \(z_j\in\mathbb R^{J-1}\), let \(u_j(z_j)\in\mathbb R^J\) be the
vector with \(j\)-th entry equal to \(1\) and remaining entries equal to
\(z_j\).

\section{SER-relax-prototype}

SER-relax-prototype is the original single-effect regression model with
uncertainty in the active LD column. The model is
\begin{align}
x\mid \beta,\gamma=j,d_j,z_j
&\sim
\Normal_J\left(\beta d_j u_j(z_j),\frac{\Rbar}{N}\right),\\
\beta&\sim \Normal(0,\sigma_0^2),\\
\gamma&\sim \Unif\{1,\ldots,J\}.
\end{align}
The relaxed LD column is induced by an inverse-Wishart model for the
unknown correlation matrix \(R\),
\[
R\sim IW_J(\nu_0\Rbar,\nu_0+J+1).
\]
For a fixed column \(j\), define \(d_j=R_{jj}\) and
\(z_j=R_{-j,j}/R_{jj}\). The induced column model is
\begin{align}
d_j&\sim
\IG\left(\frac{\nu_0+2}{2},\frac{\nu_0}{2}\right),\\
z_j\mid \omega_j
&\sim
\Normal_{J-1}\left(\mu_j,\frac{S_j}{\omega_j}\right),\\
\omega_j&\sim
\Gam\left(\frac{\nu_0+3}{2},\frac12\right),
\end{align}
using the shape--rate parameterization for Gamma distributions. Larger
\(\nu_0\) concentrates \(R\) around \(\Rbar\); smaller \(\nu_0\) allows
the active LD column to move farther away from the reference LD.

\subsection{Variational Inference}

We use the mean-field approximation
\[
q(\gamma=j,\beta,d_j,z_j,\omega_j)
=
\pi_j q_j(\beta)q_j(d_j)q_j(z_j)q_j(\omega_j),
\]
where
\[
q_j(\beta)=\Normal(m_{\beta j},v_{\beta j}),\qquad
q_j(z_j)=\Normal_{J-1}(m_{zj},V_{zj}),
\]
and
\[
q_j(\omega_j)=\Gam(\alpha_{\omega j},\beta_{\omega j}).
\]
Let
\[
\bar\beta_j=m_{\beta j},\quad
\overline{\beta^2}_j=v_{\beta j}+m_{\beta j}^2,\quad
\bar d_j=\E_j[d_j],\quad
\overline{d_j^2}=\E_j[d_j^2],\quad
\bar\omega_j=\E_j[\omega_j].
\]
Define
\begin{align}
A_j
&=
1+\tr(S_j^{-1}V_{zj})
+(m_{zj}-\mu_j)^\top S_j^{-1}(m_{zj}-\mu_j),\\
B_j
&=
x_j+(m_{zj}-\mu_j)^\top S_j^{-1}(x_{-j}-\mu_jx_j),\\
Q_j
&=
\tr(S_j^{-1}V_{zj})
+(m_{zj}-\mu_j)^\top S_j^{-1}(m_{zj}-\mu_j).
\end{align}
Then the Gaussian update for the effect size is
\begin{align}
v_{\beta j}
&=
\left(\sigma_0^{-2}+N\overline{d_j^2}A_j\right)^{-1},\\
m_{\beta j}
&=
v_{\beta j}N\bar d_jB_j.
\end{align}
Let
\[
T_{1j}=\bar\beta_j\bar d_j,\qquad
T_{2j}=\overline{\beta^2}_j\overline{d_j^2}.
\]
The Gaussian update for the relaxed LD column is
\begin{align}
V_{zj}
&=
\frac{S_j}{NT_{2j}+\bar\omega_j},\\
m_{zj}
&=
\mu_j+
\frac{NT_{1j}}{NT_{2j}+\bar\omega_j}
(x_{-j}-\mu_jx_j).
\end{align}
The Gamma factor for \(\omega_j\) is
\begin{align}
\alpha_{\omega j}
&=
\frac{\nu_0+J+2}{2},\\
\beta_{\omega j}
&=
\frac{1+Q_j}{2}.
\end{align}
The factor \(q_j(d_j)\) is non-standard. Its optimal density is
\[
q_j(d_j)
\propto
d_j^{-(\nu_0+2)/2-1}
\exp\left\{
-\frac{\nu_0}{2d_j}
-\frac{N}{2}\overline{\beta^2}_jA_jd_j^2
+N\bar\beta_jB_jd_j
\right\},
\qquad d_j>0.
\]
In implementation we work with \(r_j=\log d_j\). The unnormalized
log-density in \(r_j\), including the Jacobian, is
\[
\ell_j(r)
=
-\left(\frac{\nu_0+2}{2}+1\right)r
-\frac{\nu_0}{2e^r}
-\frac{N}{2}\overline{\beta^2}_jA_je^{2r}
+N\bar\beta_jB_je^r
+r.
\]
Numerical quadrature over \(r\) gives
\(\E_j[d_j]\), \(\E_j[d_j^2]\), \(\E_j[d_j^{-1}]\),
\(\E_j[\log d_j]\), and \(E_j[\log q_j(d_j)]\).

The coordinate probability is
\[
\pi_j=
\frac{\exp(\mathcal L_j)}
{\sum_{k=1}^J\exp(\mathcal L_k)},
\]
where the local lower bound is
\begin{align}
\mathcal L_j
&=
\log(1/J)
+E_j\log p(x\mid \beta,d_j,z_j,\gamma=j)
+E_j\log p(\beta)
+E_j\log p(d_j\mid\nu_0) \nonumber\\
&\quad
+E_j\log p(\omega_j\mid\nu_0)
+E_j\log p(z_j\mid\omega_j)
-E_j\log q_j(\beta)
-E_j\log q_j(d_j) \nonumber\\
&\quad
-E_j\log q_j(\omega_j)
-E_j\log q_j(z_j).
\end{align}
The full SER-relax-prototype ELBO is
\[
\mathcal L_{\mathrm{prototype}}
=
\sum_{j=1}^J
\pi_j\{\mathcal L_j-\log\pi_j\}.
\]

\subsection{\texorpdfstring{Learning \(\nu_0\)}{Learning nu0}}

Holding \(q\) fixed, the \(\nu_0\)-dependent part of the prototype ELBO
is
\[
M_{\mathrm{prototype}}(\nu_0)
=
\sum_{j=1}^J
\pi_j
\left\{
E_j\log p(d_j\mid\nu_0)
+E_j\log p(\omega_j\mid\nu_0)
\right\}.
\]
With
\[
a_d=\frac{\nu_0+2}{2},\quad b_d=\frac{\nu_0}{2},\quad
a_\omega=\frac{\nu_0+3}{2},\quad b_\omega=\frac12,
\]
this objective becomes
\begin{align}
M_{\mathrm{prototype}}(\nu_0)
&=
\sum_{j=1}^J\pi_j
\Big[
a_d\log b_d-\log\Gamma(a_d)
-(a_d+1)\E_j\log d_j
-b_d\E_j(d_j^{-1}) \nonumber\\
&\qquad\qquad
+a_\omega\log b_\omega-\log\Gamma(a_\omega)
+(a_\omega-1)\E_j\log\omega_j
-b_\omega\E_j\omega_j
\Big].
\end{align}
This one-dimensional objective can be optimized directly over
\(\nu_0>0\), or over \(\eta=\log\nu_0\) with Newton updates. Direct
bounded optimization is usually more stable in practice.

\subsection{Implementation Algebra}

The key identities come from the block inverse of \(\Rbar\). Because
\(\Rbar\) is a correlation matrix,
\[
\Om_{-j,-j}=S_j^{-1},\qquad
\Om_{-j,j}=-S_j^{-1}\mu_j.
\]
Therefore
\begin{align}
u_j(z_j)^\top\Om u_j(z_j)
&=
1+(z_j-\mu_j)^\top S_j^{-1}(z_j-\mu_j),\\
u_j(z_j)^\top\Om x
&=
x_j+(z_j-\mu_j)^\top S_j^{-1}(x_{-j}-\mu_jx_j).
\end{align}
These identities avoid repeated manipulation of full \(J\times J\)
columns. The prototype is nevertheless computationally heavier than the
main method because it requires quadrature for every \(q_j(d_j)\) and a
Gamma update for every \(q_j(\omega_j)\).

\section{SER-relax}

SER-relax is the simplified single-effect model. It keeps the central
idea that the active LD column can relax around \(\Rbar\), but removes
the two latent variables that made the prototype slow:
\[
d_j=1,\qquad
\omega_j=\nu_0+3.
\]
Equivalently, with
\[
\tau^2=(\nu_0+3)^{-1},\qquad \tau^{-2}=\nu_0+3,
\]
the relaxed-column prior is
\[
z_j\mid\nu_0\sim
\Normal_{J-1}(\mu_j,\tau^2S_j).
\]
The single-effect model is
\begin{align}
x\mid\beta,\gamma=j,z_j
&\sim
\Normal_J\left(\beta u_j(z_j),\frac{\Rbar}{N}\right),\\
\beta&\sim \Normal(0,\sigma_0^2),\\
\gamma&\sim\Unif\{1,\ldots,J\}.
\end{align}

\subsection{Variational Inference}

We use
\[
q(\gamma=j,\beta,z_j)=\pi_jq_j(\beta)q_j(z_j),
\]
where
\[
q_j(\beta)=\Normal(m_{\beta j},v_{\beta j}),\qquad
q_j(z_j)=\Normal_{J-1}(m_{zj},V_{zj}).
\]
For a current \(q_j(z_j)\), define
\[
\bar\beta_j=m_{\beta j},\qquad
\overline{\beta^2}_j=v_{\beta j}+m_{\beta j}^2,
\]
and the same \(A_j\), \(B_j\), and \(Q_j\) as above, but without \(d_j\)
and \(\omega_j\). The update for \(q_j(\beta)\) is
\begin{align}
v_{\beta j}
&=
\left(\sigma_0^{-2}+NA_j\right)^{-1},\\
m_{\beta j}
&=
v_{\beta j}NB_j.
\end{align}
The update for \(q_j(z_j)\) is
\begin{align}
V_{zj}
&=
\frac{S_j}{N\overline{\beta^2}_j+\tau^{-2}},\\
m_{zj}
&=
\mu_j+
\frac{N\bar\beta_j}
{N\overline{\beta^2}_j+\tau^{-2}}
(x_{-j}-\mu_jx_j).
\end{align}
The local lower bound is
\begin{align}
\mathcal L_j^{\mathrm{SER}}
&=
\log(1/J)
+E_j\log p(x\mid\beta,z_j,\gamma=j)
+E_j\log p(\beta)
+E_j\log p(z_j\mid\nu_0) \nonumber\\
&\quad
-E_j\log q_j(\beta)
-E_j\log q_j(z_j).
\end{align}
Thus
\[
\pi_j=
\frac{\exp(\mathcal L_j^{\mathrm{SER}})}
{\sum_{k=1}^J\exp(\mathcal L_k^{\mathrm{SER}})}.
\]
The expected log-likelihood is
\begin{align}
E_j\log p(x\mid\beta,z_j,\gamma=j)
&=
-\frac{J}{2}\log(2\pi)
-\frac12\log\left|\frac{\Rbar}{N}\right| \nonumber\\
&\quad
-\frac{N}{2}
\left[
x^\top\Om x
-2\bar\beta_jB_j
+\overline{\beta^2}_jA_j
\right].
\end{align}
The relaxed-column prior contribution is
\[
E_j\log p(z_j\mid\nu_0)
=
-\frac{J-1}{2}\log(2\pi)
-\frac12\log|S_j|
+\frac{J-1}{2}\log\tau^{-2}
-\frac{\tau^{-2}}{2}Q_j.
\]
The full SER-relax ELBO is
\[
\mathcal L_{\mathrm{SER}}
=
\sum_{j=1}^J
\pi_j
\{\mathcal L_j^{\mathrm{SER}}-\log\pi_j\}.
\]

\subsection{\texorpdfstring{Learning \(\nu_0\)}{Learning nu0}}

Holding \(q\) fixed, only the Gaussian prior terms for \(z_j\) depend on
\(\nu_0\). Hence, up to constants,
\[
M_{\mathrm{SER}}(\nu_0)
=
\sum_{j=1}^J
\pi_j
\left\{
\frac{J-1}{2}\log\tau^{-2}
-\frac{\tau^{-2}}{2}Q_j
\right\}.
\]
Since \(\sum_j\pi_j=1\), the closed-form update in the \(\tau^2\)
parameterization is
\[
\hat\tau^2
=
\frac{\sum_{j=1}^J\pi_jQ_j}{J-1},
\qquad
\hat\nu_0=\frac{1}{\hat\tau^2}-3.
\]
In code we bound \(\nu_0\) to avoid invalid or numerically extreme
updates. Since \(\nu_0>0\), the natural prior variance satisfies
\(\tau^2<1/3\).

\subsection{Implementation Algebra}

SER-relax can be implemented without explicitly forming all
\(S_j^{-1}\) matrices. For a vector \(r\), define
\[
C_j(r)
=
(r_{-j}-\mu_jr_j)^\top S_j^{-1}(r_{-j}-\mu_jr_j).
\]
The block inverse identity gives
\[
C_j(r)=r^\top\Om r-r_j^2.
\]
Also, because \(\Rbar_{jj}=1\), the Schur determinant identity gives
\[
|S_j|=|\Rbar|.
\]
Write
\[
z_{\mathrm{denom},j}
=
N\overline{\beta^2}_j+\tau^{-2},
\qquad
z_{\mathrm{scale},j}
=
\frac{N\bar\beta_j}{z_{\mathrm{denom},j}}.
\]
Then
\[
m_{zj}=\mu_j+z_{\mathrm{scale},j}(r_{-j}-\mu_jr_j),
\qquad
V_{zj}=S_j/z_{\mathrm{denom},j}.
\]
The quantities needed for the beta update and the local ELBO are
\begin{align}
Q_j
&=
\frac{J-1}{z_{\mathrm{denom},j}}
+z_{\mathrm{scale},j}^2C_j(r),\\
A_j&=1+Q_j,\\
B_j&=r_j+z_{\mathrm{scale},j}C_j(r).
\end{align}
Thus each coordinate only needs \(r_j\), \(C_j(r)\), the shared
\(\log|\Rbar|\), and the current scalar variational parameters.

\section{SuSiE-relax}

SuSiE-relax is the multi-effect version of SER-relax. This is the main
method. For effects \(\ell=1,\ldots,L\),
\[
\gamma_\ell\sim\Unif\{1,\ldots,J\},\qquad
\beta_\ell\sim\Normal(0,\sigma_\ell^2),
\]
and, conditional on \(\gamma_\ell=j\),
\[
z_{\ell j}\mid\nu_0
\sim
\Normal_{J-1}(\mu_j,\tau^2S_j),
\qquad
\tau^{-2}=\nu_0+3.
\]
The observation model is
\[
x\mid\{\beta_\ell,\gamma_\ell,z_{\ell,\gamma_\ell}\}_{\ell=1}^L
\sim
\Normal_J
\left(
\sum_{\ell=1}^L
\beta_\ell u_{\gamma_\ell}(z_{\ell,\gamma_\ell}),
\frac{\Rbar}{N}
\right).
\]

\subsection{Variational Inference}

We use the independent effect approximation
\[
q=
\prod_{\ell=1}^L
q_\ell(\gamma_\ell)
\prod_{j=1}^J
q_{\ell j}(\beta_\ell)q_{\ell j}(z_{\ell j}),
\]
where
\[
q_\ell(\gamma_\ell=j)=\pi_{\ell j},\qquad
q_{\ell j}(\beta_\ell)
=
\Normal(m_{\beta,\ell j},v_{\beta,\ell j}),
\]
and
\[
q_{\ell j}(z_{\ell j})
=
\Normal_{J-1}(m_{z,\ell j},V_{z,\ell j}).
\]
Define the posterior mean contribution of effect \(\ell\) by
\[
\bar\theta_\ell
=
\sum_{j=1}^J
\pi_{\ell j}\bar\beta_{\ell j}u_j(m_{z,\ell j}).
\]
When updating effect \(\ell\), use the current residual
\[
r_\ell=x-\sum_{k\ne\ell}\bar\theta_k.
\]
For each coordinate \(j\), define
\begin{align}
A_{\ell j}
&=
1+\tr(S_j^{-1}V_{z,\ell j})
+(m_{z,\ell j}-\mu_j)^\top S_j^{-1}
(m_{z,\ell j}-\mu_j),\\
B_{\ell j}(r_\ell)
&=
r_{\ell j}
+(m_{z,\ell j}-\mu_j)^\top S_j^{-1}
(r_{\ell,-j}-\mu_jr_{\ell j}),\\
Q_{\ell j}
&=
\tr(S_j^{-1}V_{z,\ell j})
+(m_{z,\ell j}-\mu_j)^\top S_j^{-1}
(m_{z,\ell j}-\mu_j).
\end{align}
Let
\[
\bar\beta_{\ell j}=m_{\beta,\ell j},\qquad
\overline{\beta^2}_{\ell j}
=
v_{\beta,\ell j}+m_{\beta,\ell j}^2.
\]
The updates are the single-effect SER-relax updates applied to the
residual \(r_\ell\):
\begin{align}
v_{\beta,\ell j}
&=
\left(\sigma_\ell^{-2}+NA_{\ell j}\right)^{-1},\\
m_{\beta,\ell j}
&=
v_{\beta,\ell j}NB_{\ell j}(r_\ell),\\
V_{z,\ell j}
&=
\frac{S_j}{N\overline{\beta^2}_{\ell j}+\tau^{-2}},\\
m_{z,\ell j}
&=
\mu_j+
\frac{N\bar\beta_{\ell j}}
{N\overline{\beta^2}_{\ell j}+\tau^{-2}}
(r_{\ell,-j}-\mu_jr_{\ell j}).
\end{align}
The coordinate weights are
\[
\pi_{\ell j}
=
\frac{
\exp(\mathcal L_{\ell j}^{\mathrm{coord}})
}{
\sum_{k=1}^J
\exp(\mathcal L_{\ell k}^{\mathrm{coord}})
},
\]
where \(\mathcal L_{\ell j}^{\mathrm{coord}}\) is the SER-relax local
lower bound computed with residual \(r_\ell\), prior variance
\(\sigma_\ell^2\), and the variational factors for effect \(\ell\).
After updating \(\pi_{\ell j}\), recompute \(\bar\theta_\ell\).

\subsection{Full ELBO}

Let
\[
C_x=
-\frac{J}{2}\log(2\pi)
-\frac12\log\left|\frac{\Rbar}{N}\right|.
\]
The expected log-likelihood is
\begin{align}
E_q\log p(x\mid\cdot)
&=
C_x
-\frac{N}{2}
\Bigg[
x^\top\Om x
-2\sum_{\ell=1}^L\bar\theta_\ell^\top\Om x \nonumber\\
&\qquad
+\sum_{\ell=1}^L\sum_{j=1}^J
\pi_{\ell j}\overline{\beta^2}_{\ell j}A_{\ell j}
+2\sum_{1\le \ell<k\le L}
\bar\theta_\ell^\top\Om\bar\theta_k
\Bigg].
\end{align}
The full SuSiE-relax ELBO is
\begin{align}
\mathcal L_{\mathrm{SuSiE-relax}}
&=
E_q\log p(x\mid\cdot) \nonumber\\
&\quad
+\sum_{\ell=1}^L\sum_{j=1}^J
\pi_{\ell j}
\Big[
\log(1/J)
+E_{\ell j}\log p(\beta_\ell)
+E_{\ell j}\log p(z_{\ell j}\mid\nu_0) \nonumber\\
&\qquad\qquad
-E_{\ell j}\log q_{\ell j}(\beta_\ell)
-E_{\ell j}\log q_{\ell j}(z_{\ell j})
\Big]
-\sum_{\ell=1}^L\sum_{j=1}^J
\pi_{\ell j}\log\pi_{\ell j}.
\end{align}
Here
\[
E_{\ell j}\log p(z_{\ell j}\mid\nu_0)
=
-\frac{J-1}{2}\log(2\pi)
-\frac12\log|S_j|
+\frac{J-1}{2}\log\tau^{-2}
-\frac{\tau^{-2}}{2}Q_{\ell j}.
\]

\subsection{\texorpdfstring{Learning \(\nu_0\)}{Learning nu0}}

The \(\nu_0\)-dependent part of the SuSiE-relax ELBO is
\[
M_{\mathrm{SuSiE-relax}}(\nu_0)
=
\sum_{\ell=1}^L\sum_{j=1}^J
\pi_{\ell j}
\left\{
\frac{J-1}{2}\log\tau^{-2}
-\frac{\tau^{-2}}{2}Q_{\ell j}
\right\}.
\]
Because \(\sum_j\pi_{\ell j}=1\) for each effect,
\[
\hat\tau^2
=
\frac{
\sum_{\ell=1}^L\sum_{j=1}^J\pi_{\ell j}Q_{\ell j}
}{L(J-1)},
\qquad
\hat\nu_0=\frac{1}{\hat\tau^2}-3.
\]
In implementation this update is bounded, for example by
\(\nu_0\in[\nu_{\min},\nu_{\max}]\). Empirically, monitoring both the
ELBO and \(\nu_0\) is important: with small reference panels, \(\nu_0\)
can move slowly and apparent duplicated effects may be transient until
the relaxation level has converged.

\subsection{Implementation Algebra}

The same scalar representation used in SER-relax applies to each
effect-specific residual \(r_\ell\). Define
\[
C_{\ell j}
=
C_j(r_\ell)
=
r_\ell^\top\Om r_\ell-r_{\ell j}^2.
\]
With
\[
z_{\mathrm{denom},\ell j}
=
N\overline{\beta^2}_{\ell j}+\tau^{-2},
\qquad
z_{\mathrm{scale},\ell j}
=
\frac{N\bar\beta_{\ell j}}
{z_{\mathrm{denom},\ell j}},
\]
we have
\begin{align}
Q_{\ell j}
&=
\frac{J-1}{z_{\mathrm{denom},\ell j}}
+z_{\mathrm{scale},\ell j}^2C_{\ell j},\\
A_{\ell j}&=1+Q_{\ell j},\\
B_{\ell j}(r_\ell)
&=
r_{\ell j}+z_{\mathrm{scale},\ell j}C_{\ell j}.
\end{align}
This is the main computational simplification: each effect update only
requires the residual \(r_\ell\), the quadratic
\(r_\ell^\top\Om r_\ell\), and coordinate-wise scalar operations.

The posterior mean effect can also be computed without explicitly
forming every \(m_{z,\ell j}\). Since
\[
u_j(m_{z,\ell j})
=
\Rbar_{\cdot j}
+z_{\mathrm{scale},\ell j}
\{r_\ell-r_{\ell j}\Rbar_{\cdot j}\},
\]
the effect mean is
\begin{align}
\bar\theta_\ell
&=
\sum_{j=1}^J
\pi_{\ell j}\bar\beta_{\ell j}
\left[
\Rbar_{\cdot j}
+z_{\mathrm{scale},\ell j}
\{r_\ell-r_{\ell j}\Rbar_{\cdot j}\}
\right] \nonumber\\
&=
\Rbar\,
\left\{
\pi_\ell\odot\bar\beta_\ell
\odot(1-z_{\mathrm{scale},\ell}\odot r_\ell)
\right\}
+
r_\ell
\sum_{j=1}^J
\pi_{\ell j}\bar\beta_{\ell j}z_{\mathrm{scale},\ell j}.
\end{align}
This is the formula used by the fast implementation. It avoids building
a \(J\times J\) matrix of relaxed columns for every effect.

\subsection{IBSS Algorithm}

The coordinate-ascent procedure is an IBSS-style algorithm:
\begin{enumerate}[leftmargin=*]
\item Initialize \(\pi_{\ell j}\), \(m_{\beta,\ell j}\),
\(v_{\beta,\ell j}\), \(m_{z,\ell j}\), \(V_{z,\ell j}\),
\(\bar\theta_\ell\), and \(\nu_0\).
\item For each effect \(\ell=1,\ldots,L\), form
\[
r_\ell=x-\sum_{k\ne\ell}\bar\theta_k.
\]
\item For all coordinates \(j\), update \(q_{\ell j}(\beta_\ell)\),
\(q_{\ell j}(z_{\ell j})\), \(A_{\ell j}\), \(B_{\ell j}\), and
\(Q_{\ell j}\) using the scalar algebra above.
\item Compute the coordinate lower bounds
\(\mathcal L_{\ell j}^{\mathrm{coord}}\), update
\(\pi_{\ell j}\) by softmax, and recompute \(\bar\theta_\ell\).
\item After sweeping through all effects, update the shared \(\nu_0\)
using the closed-form moment update, with bounds.
\item Compute the full ELBO and repeat until both the ELBO and \(\nu_0\)
have stabilized.
\end{enumerate}

\end{document}
